\section{서론}
\label{sec:intro}

\begin{figure*}[t]
  \centering
  \includegraphics[width=\linewidth]{figures/fig_architecture.pdf}
  \caption{프레임워크. 오프라인 단계에서 VQ-VAE는 보행자 궤적으로부터 $K=16$개의 보행 행동으로 이루어진 코드북과 사전 네트워크를
  학습한다. 각 상황에서 코드북은 후보 경로들로 디코딩되고, 안전하지 않은 후보는 제거되며, 각 안전한 후보는 다섯 가지 스타일 특징으로
  기술된다. 소수의 질문 상황에서 EUBO가 안전한 후보 한 쌍을 고르고, 사용자는 어느 쪽을 선호하는지 답하며, 선호에 대한 베이지안
  Bradley--Terry 사후분포가 갱신된다. 새로운 상황에서는 추가 질문 없이, 로봇이 사후 평균 $\hat\theta$ 하에서 가장 좋은 안전한 후보를
  실행한다. Fig.~\ref{fig:overview}는 실제 상황에서의 각 단계를 보여준다. \texttt{utils/figure\_architecture.py}로 생성하였다.}
  \label{fig:architecture}
\end{figure*}

로봇이 보도, 역, 상점가를 사람들과 공유하는 일이 점점 늘어나고 있으며, 이러한 로봇이 군중 속을 어떻게 이동하는지는 로봇이 서비스하는 사람,
즉 배달 로봇을 보내거나 휠체어를 타거나 안내 로봇과 나란히 걷는 사용자에게 중요하다. 그러나 ``사회적으로 순응하는''
내비게이션은 하나의 단일한 행동이 아니다. 마주 오는 보행자를 만났을 때, 어떤 사용자는 로봇이 속도를 줄이고 상대를 먼저 지나가게 하기를 원하고,
다른 사용자는 속도를 유지하기를 원한다. 어떤 사용자는 로봇이 사람들의 뒤로 지나가기를 선호하고, 다른 사용자는 앞으로 지나가기를 선호한다.
어떤 사용자는 무리를 크게 우회하는 것을 감수하지만, 다른 사용자는 가장 짧은 안전한 경로를 원한다. 사람들 자신도 타인을 얼마나 가까이 지나치는지,
그리고 얼마나 쉽게 길을 양보하는지에서 서로 다르므로~\citep{hall1966hidden}, 바람직한 군중 행동에 대한 단일한 개념으로는 많은 사용자를
만족시키지 못한다.

그럼에도 군중 내비게이션 방법들은 하나의 고정된 목적함수를 최적화한다. 모델 기반 플래너는 수작업으로 조정된 사회적 비용을 인코딩하고, 강화학습
정책은 설계된 보상을 최대화하며, 모방 학습 또는 생성형 플래너는 학습 데이터의 평균적인 보행자를 재현한다
\citep{helbing1995social,vandenberg2011reciprocal,chen2019crowd,kretzschmar2016socially,mavrogiannis2023core}. 이들 중 어느 것이든 개인에게 맞추려면
재조정이나 재학습이 필요하다. 선호 기반 학습은 이에 대한 직접적인 해법을 제공한다. 즉, 보상을 명시하는 대신 사용자가 두 행동을 비교하면 로봇이
사용자가 원하는 바를 추론한다~\citep{christiano2017deep,sadigh2017active}. 그러나 로보틱스에서는 각 비교가 사용자의 시간과 주의를 소모하므로,
소수의 응답만으로 선호를 식별해야 하며, 비교되는 궤적은 사람에게 의미 있는 것이어야 한다. 즉, 안전하고, 그럴듯하며, 사용자가 중요하게 여기는
측면에서 서로 달라야 한다.

우리의 핵심 아이디어는 로봇이 \emph{무엇을} 할 수 있는지와, 그 행동들 중 특정 사용자가 \emph{어떤 것을} 원하는지를 분리하는 것이다(Fig.~\ref{fig:architecture}).
전자를 위해, 보행자 궤적으로 학습한 VQ-VAE~\citep{vandenoord2017neural}를 이용하여 인간 보행 행동의 이산 코드북을 학습한다. 주어진 상황에서
코드북을 디코딩하면, 구성상 인간과 유사하면서 스타일이 서로 다른 -- 왼쪽 또는 오른쪽으로, 앞 또는 뒤로 지나가거나, 속도를 줄이거나 유지하는 --
소수의 후보 경로 집합이 얻어지며, 안전 필터는 누군가에게 지나치게 가까이 접근하는 후보를 제거한다. 후자를 위해, 각 안전한 후보를 해석 가능한
다섯 가지 스타일 특징으로 기술하고, 사용자의 선호를 이 특징들에 대한 선형 효용으로 모델링하며, 이를 A/B 응답으로부터 베이지안 Bradley--Terry
모델~\citep{bradley1952rank}로 학습한다. 로봇은 질문(query)을 마친 후 각 상황에서 하나의 경로만 실행하므로 적절한 목적은 단순 regret(simple regret)이며,
이에 대해 한 단계 베이즈 최적인 최선 선택지의 기대 효용(EUBO)~\citep{astudillo2023qeubo}에 따라 각 질문을 선택한다.

데이터 사전분포만으로는 개인의 선호를 포착하지 못한다. Grand Central과 ETH/UCY의 학습에 쓰지 않은(held-out) 상황에서, 항상 가장 인간과 유사한
안전 경로를 선택하면 시뮬레이션된 사용자에 대해 정규화된 단순 regret이 0.45--0.47로 남으며, 이는 무작위 안전 경로(0.49)와 비슷한 수준이다. 다섯 번의 EUBO
질문은 이를 0.09--0.12로 줄이며, 네 가지 설정 모두에서 무작위 질문보다, 그리고 세 가지 설정에서 능동 질문 기준선보다 낮은 regret을 보이며,
정보 이득 질문만큼 낮은 regret을 달성한다. 무작위 질문 대비 이점은 노이즈가 있는, 비일관적인, 그리고 비선형 시뮬레이션 사용자에 대해서도 유지된다.
현재 선호 미세조정의 상당 부분을 뒷받침하는 추정기인 Direct Preference Optimization~\citep{rafailov2023direct}과의 비교는
이득이 어디에서 오는지를 보여준다. 동일한 질문이 주어지면 DPO와 Bradley--Terry 사후분포는 비슷한 성능을 보이지만, DPO에는 질문을 선택하는 데
사용할 사후분포가 없다. 후보 집합 또한 중요하다. 공통 Grand Central 장면에서 flow matching 기반
군중 내비게이션 플래너~\citep{singha2026crowdfm}의 후보들은 코드북보다 훨씬 좁은 범위의 행동만을 포괄하며, 이는 어떠한 학습된 선호로도 달성할 수 있는 성능의
한계를 정한다(다섯 번의 질문 후 union regret 0.35 대 0.11). 마지막으로, 폐루프 리플레이에서 학습된 선호를 따르면 충돌률의 유의미한 변화 없이
사용자의 스텝 regret이 0.10--0.14로 유지된다.

\paragraph{기여.}
\begin{itemize}
  \item 보행자 궤적으로부터 VQ-VAE로 학습한 보행 행동 코드북. 이는 하나의 상황을 안전 필터가 검사할 수 있는, 다양하고 인간과 유사한
        소수의 후보 경로 집합으로 변환한다(Sec.~\ref{sec:candidates}).
  \item 해석 가능한 다섯 가지 스타일 특징에 대한 선형 선호. 이는 A/B 응답으로부터 베이지안 Bradley--Terry 모델로 학습되며,
        질문은 단순 regret을 최소화하도록 EUBO로 선택된다(Sec.~\ref{sec:preference}).
  \item 학습에 쓰지 않은 상황에 대한 단순 regret 평가 프로토콜(Sec.~\ref{sec:regret}), 그리고 질문 선택(정보 이득 및 모델 밖의 사용자 포함), DPO 추정기, flow matching 기반
        후보 생성기, 폐루프 실행을 비교하는 Grand Central 및 ETH/UCY 실험(Sec.~\ref{sec:experiments}).
\end{itemize}
